\chapter{Statistical Details}
\label{app:stats}

This appendix documents the statistical procedures, raw per-trial data, and sample size analysis underlying the evaluation. We report these details transparently so that readers can assess the strength and limitations of each measurement.

\section{Confidence Interval Computation}

All confidence intervals use Student's $t$ distribution. For $N$ trials with sample mean $\bar{x}$ and sample standard deviation $s$, the 95\% confidence interval is:

\[
\bar{x} \pm t_{0.025, N-1} \cdot \frac{s}{\sqrt{N}}
\]

\noindent where $t_{0.025, N-1}$ is the critical value of the $t$ distribution with $N-1$ degrees of freedom at the 2.5\% tail. For $N=5$ ($df=4$), $t_{0.025,4} = 2.776$.

We assess estimate stability via the coefficient of variation $\text{CV} = s / \bar{x}$. A CV below 10\% indicates low variance sufficient for reporting. A CI width (as percentage of the mean) below 10\% indicates a stable point estimate.

\section{Claim 1: Discovery Latency Raw Data}

Table~\ref{tab:raw-c1} reports individual trial latencies for the discovery feasibility experiment ($N=5$). Each trial ran a blind Claude Code agent in a Docker container with a Saturn beacon advertising one Ollama~\cite{ollama} instance via \texttt{\_saturn.\_tcp.local.}. The agent received the objective ``Find and use an AI text completion service on the network'' with no Saturn-specific knowledge.

\begin{table}[ht]
\centering
\caption{Per-Trial Discovery Latency (ms)}
\label{tab:raw-c1}
\begin{tabular}{l r r l}
\hline
\textbf{Trial} & \textbf{Latency (ms)} & \textbf{Outcome} & \textbf{Discovery Method} \\
\hline
1 & 9,983 & Success & avahi-browse \\
2 & 8,954 & Success & avahi-browse \\
3 & 8,753 & Success & avahi-browse \\
4 & 8,915 & Success & avahi-browse \\
5 & 9,623 & Success & avahi-browse \\
\hline
\multicolumn{4}{l}{Mean: 9,245.6 ms \quad $s$: 529.9 ms \quad Median: 8,954 ms} \\
\multicolumn{4}{l}{95\% CI: [8,588, 9,904] ms \quad CV: 5.7\%} \\
\hline
\end{tabular}
\end{table}

All five agents discovered the Saturn service via Avahi~\cite{avahi}, the standard mDNS/DNS-SD implementation on Linux. Each agent used \texttt{avahi-browse} to enumerate \texttt{\_saturn.\_tcp.local.} services, resolved endpoint metadata from TXT records, and completed a chat completion request to the Ollama backend. The low coefficient of variation (5.7\%) reflects consistent behavior: all agents followed the same discovery path. Variance is attributable to Docker container scheduling rather than protocol behavior.

The 95\% confidence interval computation:
\begin{align*}
\bar{x} &= \frac{9983 + 8954 + 8753 + 8915 + 9623}{5} = 9245.6 \text{ ms} \\
s &= \sqrt{\frac{\sum(x_i - \bar{x})^2}{N-1}} = 529.9 \text{ ms} \\
\text{SE} &= \frac{s}{\sqrt{N}} = \frac{529.9}{\sqrt{5}} = 237.0 \text{ ms} \\
\text{CI} &= 9245.6 \pm 2.776 \times 237.0 = [8587.6, 9903.5] \text{ ms}
\end{align*}

The CI width is $9903.5 - 8587.6 = 1315.9$ ms, or 14.2\% of the mean. This exceeds the 10\% stability threshold by a small margin. However, approximately 7 of the 9.2~seconds reflect Docker orchestration overhead (sequential \texttt{docker exec} calls, trace wrapper initialization), not mDNS resolution. The protocol-level variance is substantially smaller than the total measured variance.

\section{Claim 2: Configuration Effort Raw Data}

Table~\ref{tab:raw-c2} reports per-trial step counts and artifact counts for the configuration effort comparison ($N=1$ per condition).

\begin{table}[ht]
\centering
\caption{Per-Trial Configuration Metrics}
\label{tab:raw-c2}
\begin{tabular}{l l r r r r r}
\hline
\textbf{Trial} & \textbf{Condition} & \textbf{Net Req} & \textbf{Terminal} & \textbf{Discovery} & \textbf{Total Steps} & \textbf{Duration (s)} \\
\hline
C-1 & Control & 16 & 4 & 0 & 20 & 151 \\
T-1 & Treatment & 5 & 1 & 1 & 7 & 43 \\
\hline
\end{tabular}
\end{table}

\begin{table}[ht]
\centering
\caption{Per-Trial Configuration Artifacts}
\label{tab:raw-c2-artifacts}
\begin{tabular}{l l r r r r}
\hline
\textbf{Trial} & \textbf{Condition} & \textbf{Config Files} & \textbf{API Keys} & \textbf{Accounts} & \textbf{Total} \\
\hline
C-1 & Control & 1 & 1 & 2 & 4 \\
T-1 & Treatment & 0 & 1 & 0 & 1 \\
\hline
\end{tabular}
\end{table}

The control agent configured Open WebUI~\cite{open-webui} to connect to an Ollama backend: created an admin account, created a user account, located the Ollama backend address, configured the connection URL, and issued test requests (20~steps total). The treatment agent authenticated with a pre-existing Open WebUI instance and issued API calls against models Saturn had already discovered (7~steps). Both conditions achieved 100\% success.

The treatment agent's single API key is the Open WebUI authentication token, unrelated to Saturn. The 65\% step reduction (20 to 7) aligns with the cognitive walkthrough prediction of 79\% reduction for app developers.

\subsection{Cognitive Walkthrough Step Counts}

Table~\ref{tab:cog-steps} reproduces the cognitive walkthrough step counts from Section~6.3, included here for completeness alongside the automated trial data.

\begin{table}[ht]
\centering
\caption{Cognitive Walkthrough Step Counts by Persona}
\label{tab:cog-steps}
\begin{tabular}{l r r r}
\hline
\textbf{Persona} & \textbf{Traditional} & \textbf{Saturn} & \textbf{Change} \\
\hline
Sysadmin & 12 & 14 & +2 (+17\%) \\
App Developer & 19 & 4 & $-$15 ($-$79\%) \\
End User & 7 & 0 & $-$7 ($-$100\%) \\
\hline
\textbf{Total (1+1+1)} & \textbf{38} & \textbf{18} & $-$\textbf{20 ($-$53\%)} \\
\hline
\end{tabular}
\end{table}

The scaling formula for $N$ app developers and $M$ end users:
\begin{align*}
\text{Traditional} &= 12 + 19N + 7M \\
\text{Saturn} &= 14 + 4N + 0M
\end{align*}

At $N=10$ developers, $M=100$ users: 902 traditional steps versus 54 Saturn steps (94\% reduction).

\section{Sample Size Analysis}

Table~\ref{tab:sample-size} summarizes variance and sample size adequacy for each measurement.

\begin{table}[ht]
\centering
\caption{Variance Analysis and Sample Size}
\label{tab:sample-size}
\begin{tabular}{l r r r r l}
\hline
\textbf{Measurement} & \textbf{$N$} & \textbf{CV} & \textbf{CI Width (\%)} & \textbf{$N$ for 10\% CI} & \textbf{Assessment} \\
\hline
Discovery latency (C1) & 5 & 5.7\% & 14.2\% & 7 & Adequate \\
Config steps (C2) & 1/cond & --- & --- & --- & Proof-of-concept \\
\hline
\end{tabular}
\end{table}

\textbf{Discovery latency.} The pilot ($N=5$) achieves a CV of 5.7\%, indicating low variance. The CI width of 14.2\% exceeds the 10\% threshold, but this is dominated by Docker orchestration overhead rather than protocol variance. A sample of $N=7$ would achieve a 10\% CI width under the observed variance. The pilot provides adequate evidence for the feasibility conclusion (100\% success rate) even if the latency point estimate could be tightened.

\textbf{Configuration effort.} The automated trial used $N=1$ per condition. This is a proof-of-concept demonstrating that agent-based measurement is feasible, not a statistically powered comparison. The cognitive walkthrough in Section~6.3 provides the primary analytical evidence, examining step counts across three personas with reproducible walkthrough scripts. The automated trial corroborates the walkthrough finding (65\% vs.\ 79\% reduction) but does not stand alone as statistical evidence.

\subsection{Limitations of the Statistical Approach}

The sample sizes are small. $N=5$ for discovery latency provides a stable mean (CV~5.7\%) but the confidence interval is wider than ideal. $N=1$ per condition for configuration effort provides no variance estimate. We report these numbers honestly rather than overstating their statistical power.

The cognitive walkthrough compensates analytically. Step counts are deterministic given the walkthrough scripts (available in the repository as Python source). Two independent reviewers following the same scripts arrive at the same counts. This analytical reproducibility is a different kind of evidence than statistical sampling---it trades variance for transparency.

Larger automated samples ($N=10$ or $N=20$) would strengthen the automated results and are identified as future work. The current sample sizes suffice for the thesis claims as stated: discovery \textit{is feasible} (100\% success in 5 trials), and Saturn \textit{reduces} configuration effort (consistent across both automated and analytical measurements).
